% Generated by `taskgen course examples` from ecc.naf (seed 1).
% Do not edit: change examples.yaml and regenerate.
\begin{example}[NAF и число операций]\label{ex:ecdlp-naf}
Найти несмежную форму (NAF, тернарное разложение с цифрами $-1, 0, 1$) числа $n = 442$. Сколько удвоений и сложений (вычитаний) точек требуется для вычисления $nP$ на эллиптической кривой по бинарному разложению и по NAF?

\par\noindent\textit{Решение.} Двоичная запись: $n = 110111010_2$ (длина $9$, единиц $6$). Цифры NAF справа налево: если $n_i$ чётно, $z_i = 0$; иначе $z_i = 2 - (n_i \bmod 4) \in \{1, -1\}$; $n_{i+1} = (n_i - z_i)/2$.\[\begin{array}{c|r|c|r|r}i & n_i & n_i \bmod 4 & z_i & n_{i+1} \\ \hline 0 & 442 & 2 & 0 & 221 \\ 1 & 221 & 1 & 1 & 110 \\ 2 & 110 & 2 & 0 & 55 \\ 3 & 55 & 3 & -1 & 28 \\ 4 & 28 & 0 & 0 & 14 \\ 5 & 14 & 2 & 0 & 7 \\ 6 & 7 & 3 & -1 & 4 \\ 7 & 4 & 0 & 0 & 2 \\ 8 & 2 & 2 & 0 & 1 \\ 9 & 1 & 1 & 1 & 0 \\ \end{array}\]$\mathrm{NAF}(442) = (100\bar 100\bar 1010)$, т.\,е. $442 = 2^{9} - 2^{6} - 2^{3} + 2^{1}$ (длина $10$, ненулевых цифр $4$; соседних ненулевых цифр нет). NAF на одну цифру длиннее двоичной записи (старшая цифра $1$ стоит в разряде $2^{9} > n$).\par Начиная с $R = P$, на каждую следующую цифру выполняется одно удвоение и, если цифра ненулевая, одно сложение ($+P$) или вычитание ($-P$; $-P$ вычисляется бесплатно):\[\begin{array}{l|c|c|c} & \text{удвоений} & \text{сложений} & \text{всего} \\ \hline \text{бинарное} & 8 & 5 & 13 \\ \text{NAF} & 9 & 3 & 12 \\ \end{array}\]NAF экономит $1$ операцию. Для сравнения: $\log_2 n \approx 8{,}79$, оценки в среднем $1{,}5\log_2 n \approx 13{,}2$ (бинарное) и $\tfrac43\log_2 n \approx 11{,}7$ (NAF).\[\mathrm{NAF}(442) = (100\bar 100\bar 1010);\quad 13 \text{ и } 12 \text{ операций}.\]
\end{example}
